\section{Evaluation Design}\label{sec:eval}

Because this is a measurement study rather than a system, our evaluation asks
whether the method produces \emph{trustworthy, interpretable} evidence of
app-store censorship. We evaluate along the axes below; Section~\ref{sec:results}
reports the outcomes.

\subsection{Coverage}

We report the breadth of the measurement: the number of apps and countries
successfully queried, the fraction of the app $\times$ country matrix for which we
obtained a definitive availability signal, and the cadence and completeness of the
weekly time series. Coverage is the denominator against which every other claim is
made, so we state explicitly which storefronts or apps we could not measure and
why (e.g., endpoints that rate-limited us or apps with ambiguous identifiers).

\subsection{Validation Against Ground Truth}

The central test of the method is whether it independently reproduces removals
that are already documented. We compare our detections against (a) Apple and
Google transparency reports, (b) GreatFire AppleCensorship
records~\cite{applecensorship}, and (c) news-reported events such as the July 2024
Russian VPN removals~\cite{techcrunch2024russia,moscowtimes2024russia} and the
2017 China VPN removals~\cite{cnbc2017china}. If our pipeline surfaces these known
events from raw availability data, that validates the measurement; disagreements
are analyzed as either method limitations or genuinely new findings.

\subsection{Attribution: Separating Censorship From Coincidence}

The hardest evaluation question is the one that motivates the project: a missing
app does not by itself prove government censorship, since the cause could be a
developer decision, a licensing restriction, a sanctions regime, or a
policy-violation removal. We do not claim to resolve intent perfectly; instead we
build a defensible attribution procedure and evaluate its discriminating power:
\begin{itemize}
  \item \textbf{Control apps.} A neutral app (e.g., a weather app) missing in a
  country is a signal of a benign cause; a VPN missing in the same country is far
  more suspicious. The gap between control and sensitive categories bounds the
  false-positive rate.
  \item \textbf{Category-level patterns.} A single missing VPN could be
  coincidence; \emph{all} VPNs disappearing in one country over one week is
  consistent with a policy action, not chance.
  \item \textbf{Triangulation.} We corroborate suspected censorship events against
  transparency reports and news, and report a confidence label (documented /
  strongly inferred / ambiguous) for each.
\end{itemize}
We quantify how cleanly these heuristics separate our sensitive basket from the
controls (Section~\ref{sec:results}), and we are explicit about what the data can
and cannot establish. This honest treatment of attribution is itself a
contribution, since it directly addresses a known weakness of availability-based
censorship measurement.

\subsection{Correlation With External Indices}

Finally, we test whether per-country censorship intensity (missing sensitive apps,
normalized by controls) correlates with independent repression measures such as
\emph{Freedom on the Net}. A strong positive correlation would be external
evidence that our metric captures real censorship rather than measurement noise.
